\section{Related Work}

\subsection{Temporal Reasoning Benchmarks}
Recent years have seen a proliferation of temporal reasoning benchmarks for LLMs.
\textsc{TimE}~\citep{time2025} provides 38,522 questions across 11 subtasks at 3 difficulty levels, evaluating LLMs' ability to understand temporal expressions and reason about event sequences.
\textsc{ETRQA}~\citep{etrqa2024} consolidates event temporal reasoning datasets for reading comprehension.
\textsc{Test of Time}~\citep{testoftime2025} generates synthetic temporal reasoning tasks.
\textsc{MusTQ}~\citep{mustq2024} evaluates multilingual temporal KGQA across 10 languages.
However, all treat time as a \emph{single} dimension: none systematically distinguish event occurrence time from information recording time, a distinction critical for memory systems.

\subsection{Memory System Benchmarks}
The emergence of memory-augmented LLM agents has driven benchmark development.
LoCoMo~\citep{locomo2024} (ACL 2024) evaluates long conversational memory with 1,986 questions across 5 types, including adversarial and temporal categories.
LongMemEval~\citep{longmemeval2024} provides 4,500 questions testing 8 memory systems on single-hop, multi-hop, and temporal reasoning.
MemoryAgentBench~\citep{memoryagentbench2025} focuses on multi-dataset agent memory evaluation.
While these benchmarks include temporal questions, they do not systematically test \emph{dual-timestamp} reasoning, the ability to distinguish ``when this happened'' from ``when I learned about it.''

\subsection{Bitemporal Data Modeling}
The distinction between event time (valid time) and record time (transaction time) is foundational in temporal database research~\citep{snodgrass1999tsql2,jensen1999temporal}.
Bitemporal data models track both dimensions, enabling queries like ``what did the system know at time $T_1$ about events that occurred at time $T_2$?''
This concept has been extensively studied in database systems but has not been applied to evaluate LLM memory systems.
BiTempQA bridges this gap by bringing bitemporal reasoning into the memory system evaluation paradigm.

\subsection{Positioning}
Table~\ref{tab:benchmark_comparison} positions BiTempQA relative to existing work.
BiTempQA is the \emph{only} benchmark that (1) evaluates dual-timestamp reasoning and (2) tests multiple memory system backends.

\begin{table}[t]
\centering
\resizebox{\columnwidth}{!}{%
\begin{tabular}{lcccccc}
\toprule
\textbf{Benchmark} & \textbf{Year} & \textbf{Temp.} & \textbf{Dual TS} & \textbf{Mem} & \textbf{Lang} & \textbf{Scale} \\
\midrule
LoCoMo & 2024 & Partial & \xmark & \cmark & EN & 1,986 \\
LongMemEval & 2024 & \xmark & \xmark & \cmark & EN & 4,500 \\
TimE & 2025 & \cmark & \xmark & \xmark & EN & 38,522 \\
ETRQA & 2025 & \cmark & \xmark & \xmark & EN & Consol. \\
Test of Time & 2025 & \cmark & \xmark & \xmark & EN & Novel \\
MusTQ & 2024 & \cmark & \xmark & \xmark & Multi & 15K \\
\midrule
\textbf{BiTempQA} & \textbf{2026} & \cmark & \cmark & \cmark & ZH & 1,536 \\
\bottomrule
\end{tabular}%
}
\caption{Comparison with related benchmarks. \textbf{Dual TS}: evaluates dual-timestamp reasoning. \textbf{Mem}: tests memory backends (not just LLMs).}
\label{tab:benchmark_comparison}
\end{table}
